% Introduction
% viscode_shared_subspace_probe — D028/D029/D030 framing
% Generated: 2026-04-11

\section{Introduction}
\label{sec:intro}

% P1: visual code generation booming, all behavioral eval
The generation of visual content through code is a rapidly growing capability of large language models.
Rather than producing raster images directly, these models synthesize programs in domain-specific formats---SVG for web graphics, TikZ for scientific figures, Asymptote for geometric illustrations---that compile into precise, scalable visual output.
Specialized systems such as VisCoder2 \cite{viscoder2}, LLM4SVG \cite{llm4svg2025}, and Chat2SVG \cite{chat2svg2025} have demonstrated increasingly sophisticated visual code generation, while general-purpose code LLMs like the Qwen2.5-Coder family \cite{qwen25coder} handle multiple such formats within a single model.
However, all existing evaluations of these models are \emph{purely behavioral}: models are judged by execution pass rates, VLM-judge scores, or CLIP similarity, with no analysis of what visual structure the models internally encode.

% P2: natural question
This leaves a fundamental question unanswered: \emph{do models that generate visual code in multiple formats develop format-agnostic internal representations of visual content, or are their hidden states predominantly format-specific?}
If different visual code formats converge to similar representations for semantically related visual programs, it would suggest that these models have abstracted visual semantics beyond surface syntax.
Conversely, if representations are entirely format-specific---as one might expect given the radically different tokenization of SVG (XML markup), TikZ (\LaTeX{} macros), and Asymptote (imperative C-like code)---it would imply that multi-format capability arises from disjoint pathways rather than shared computation.

% P3: our contribution
Addressing this question is methodologically non-trivial.
Unlike probing for discrete linguistic properties \cite{hewitt-manning-2019-structural, belinkov-2022-probing}, measuring representational similarity across output formats requires comparing hidden-state geometries for semantically related but syntactically distinct programs, with careful attention to small sample sizes and syntactic confounds.
In this work, we introduce a \textbf{training-free probing framework} for measuring cross-format representational similarity in visual code LLMs.
The framework operates entirely on cached hidden-state representations extracted via teacher-forcing, combining linear Centered Kernel Alignment (CKA; \citealt{kornblith2019similarity}) with bootstrap confidence intervals, permutation null testing, and a Procrustes distance robustness check.
We demonstrate this framework through a case study on the Qwen2.5 model family.

% P4-P5: compressed — findings already in Abstract
Our case study on three Qwen2.5-based 7B models finds statistically significant but weak cross-format CKA (${\sim}3\times$ the permutation null), lexically confounded at pair-level (BPE token Jaccard $\rho \approx 0.7$), with residual layer-dependent structure.
The observed similarity is of \textbf{uncertain origin}, and we present these results as a methodological template.
Our contributions are:

\begin{enumerate}
    \item \textbf{A training-free probing framework} for measuring cross-format representational similarity in LLMs, combining CKA with bootstrap confidence intervals, permutation null testing (5{,}000 permutations), power simulation, and Procrustes distance robustness check.

    \item \textbf{Five empirical findings} from a Qwen2.5 case study---the first representation-level analysis of multi-format visual code LLMs, including a pre-registered negative control that establishes the measured cross-format signal reflects code-general rather than visual-specific structure---demonstrating both the utility of the framework and the complexity of interpreting cross-format similarity signals.

    \item \textbf{An open-source analysis pipeline} covering data generation, hidden-state extraction, CKA computation, permutation testing, and visualization, designed for replication across model families beyond the single-family case study presented here.
\end{enumerate}

